\chapter{}

Ghìsola era stata mandata via da Poggio a' Meli, con astuta precauzione, da Domenico; che, vedendo il contegno poco sicuro di lei, non volle trovarsi in impicci.

Ella era andata a Poggio a' Meli a dodici anni ed era tornata a Radda a diciassette.

Conosceva quasi soltanto di nome gli altri parenti e non aveva più veduto le due sorelle, che non le erano affezionate, perchè non vivevano insieme; ma andarono a prenderla alla diligenza, mettendosi le scarpe nuove e gli scialletti delle feste.

Ella portò loro due anelli d'oro falso, per regalo. La baciarono e poi si trovarono tutte e due impacciate. Non sapevano se la tenevano nel mezzo; e, camminando, cambiavano sempre di posto. La minore, anzi, si mise dietro; e, quando Ghìsola la chiamò con sé, invece andò lungo la proda sul margine erboso della strada; riabbassando la testa tutte le volte che Ghìsola si voltava a lei, perchè non voleva far vedere che la guardava. Anche la sorella più grande parlò poco, anzi non disse niente.

Quando giunsero a casa, dove l'aspettavano i genitori, Ghìsola si mise a piangere. Ma, poi, fecero un bel pranzo, mangiando un coniglio fritto e due galline in padella; due galline che avrebbero dovuto campare, perchè avevano le ovaie grasse e piene. Il pane era stato sfornato la mattina stessa.

Borio di Sandro, un vedovo amico della famiglia che aiutava anche con il denaro, aveva portato un fiasco del suo buon vino. E, il primo giorno, quella mezza sbornia mise tutti d'accordo.

Ma Ghìsola non se la sentiva di faticare come le sorelle, che la chiamavano tra sè la «signorina delicata». Non voleva saperne di starci insieme; e, quando le era possibile, andava nel campo sola. Non le volevano male, ma lei trovava sempre modo di smetter subito qualunque discorso che volessero incominciare. Anche alla messa andava sola; e ripensava a Poggio a' Meli. Già tornare a Radda era stato un dispiacere; e Borio soltanto lo capiva. Ella gli diceva sempre che non ci sarebbe rimasta a costo di farsi ammazzare!

Un anno dopo, la sera di una solenne festa religiosa, egli l'aveva accompagnata alla processione su dentro il paese.

Era stata una processione con i contadini dei dintorni dietro ad una piccola croce, a coppie, con i loro cappelli in mano. Le ragazze, tutte insieme dopo, cantavano leggendo in un libro tenuto aperto con ambedue le mani, sempre a testa bassa, come quando si va incontro a un vento impetuoso. Poi un altra croce, grande e nera, polverosa, con una corona di spine e con i flagelli di corda pendenti. Poi il prete.

Il vedovo ricondusse a casa Ghìsola che non aveva mai voluto dare retta a nessun giovinotto, perchè si teneva molto da più di tutti.

Scesero per una strada ripida, sempre più buia, che porta fuori del paese; accanto alle file dei cipressi folti, entrando poi nei campi. Percorsero un sentiero scosceso, a metà di un grande poggio nano e coperto di querci alte.

Ghìsola, a cui Borio piaceva molto, camminava un passo innanzi, un poco triste come succede sovente dopo l'allegria insolita e quasi involontaria di una festa.

«Perchè ella non mi guarda più?»

Gettò via il sigaro che ora gli faceva male e gli aumentava la confusione. Erano soli! Tutta l'altra gente non si sa dove fosse scomparsa! È vero che qualche volta egli udiva, prima di lei, rumore di passi; ma poi il calpisticcio si allontanava.

Pareva che Ghìsola volesse farsi sempre più piccola, camminando quasi senza vedere; e se non ci fosse stato Borio, a cui stava vicino ascoltandolo respirare, sarebbe andata a battere in qualche proda.

Di quando in quando, inciampava; le sue gambe parevano intirizzite e così lunghe che ad ogni passo la facevano rintronare tutta. E allora pensò di fermarsi. Credeva d'aver bevuto troppo; e si sentiva portar via la testa; senza avvedersene, sospirava sollevando lungamente lo stomaco.

L'oscurità, con la luna palpitante sotto un velo di nuvole, empiva ogni parte di ombre fievoli e trasparenti. Allora egli la prese per mano, ed ella lasciò fare: gli pareva che Ghìsola fosse doventata un essere debole, quasi buffa. Ma capì. La baciò; ed ella si discostò, trasalendo. La baciò ancora, guardando dopo fissamente la sua nuca e il suo dorso solcato tra le spalle. Ma, forse, non sarebbe riuscito a baciarla un'altra volta! E siccome non si voltava a dietro, le cinse la vita con il braccio.

Stava zitta! Ella aveva paura di parlare, quanto dell'ombre di quei cipressi: le quali, all'improvviso, subito fuori del paese, attraversavano la strada, risalendo come se fossero vive, con la cima su per il muro dalla parte opposta.

Ad un tratto si sedette a metà del viottolo sopra una pietra, nascondendo la faccia con lo scialletto caduto giù dai capelli; e, sopra, le mani: mani che parevano di ferro, come le punte del forcone.

Egli, volendole parlare, pur non sapendo come, dovette abbassarsi tutto. Non gli pareva di essere accanto a quella Ghìsola che conosceva da tanto tempo e che era con lui anche poco fa. Ella strinse le gambe l'una contro l'altra, così insieme che somigliavano ad un aratro voltato in sù.

Allora Borio, dopo una lotta silenziosa, con le mani, potè dire, sentendo già il rimorso, senza nessuna voluttà:

— Ti dico di sì.... ti dico di sì....

Le loro dita, sudate, si sguisciavano; egli aveva voglia di storcergliele: si guardavano come quando si sta per leticare, perchè ormai era impossibile smettere.

Ella allontanò le gambe. Poi pianse.
